\documentclass[11pt]{article}
\usepackage[utf8]{inputenc}
\usepackage[T1]{fontenc}
\usepackage{lmodern}
\usepackage{amsmath,amssymb,amsthm}
\usepackage[margin=1in]{geometry}
\usepackage{xcolor}
\usepackage{hyperref}
\hypersetup{colorlinks=true,linkcolor=blue!50!black,urlcolor=blue!50!black}

\newtheorem{question}{Question}
\newcommand{\Q}{\mathbb{Q}}
\newcommand{\F}{\mathbb{F}}
\newcommand{\Z}{\mathbb{Z}}

\title{\textbf{A transcendence question for a catalytic $q$-difference series}\\[0.2em]
\large arising from the geodesic growth of a planar reflection group (OEIS A396406)}
\author{Vico Bonfioli\thanks{Independent researcher. \texttt{vicobonfioli@gmail.com}. Code and data:
\url{https://github.com/ElVec1o/pythagorean-reflection-sequence}.}}
\date{June 2026}

\begin{document}
\maketitle

\noindent\textbf{Summary.} A natural integer power series $F(q)\in\Z[[q]]$ of radius $1$, defined by an
explicit linear $q$-difference equation with the dilation $t\mapsto q^2t$, appears to be transcendental
over $\Q(q)$, and the reduction modulo every prime $p$ appears \emph{maximally non-$p$-automatic}.
All the standard transcendence tools we have tried provably cap out at a precise obstruction (a
``natural-boundary'' degree balance). \textbf{We are looking for a transcendence criterion, or a pointer to
the right machinery or person, for this specific $q$-difference specialization.}

\section*{The object}

Let $W$ be the group generated by the three side-reflections of a right triangle in the plane with
unequal rational legs. Its geodesic growth series $U(x)=\sum_{n\ge0}u_n x^n$ is OEIS
\href{https://oeis.org/A396406}{A396406} ($u_n=$ number of group elements at word-length $n$). A
catalytic (lamplighter) transfer for the underlying length statistic, with $q=x^2$ and a catalytic
variable $t$, yields a bivariate generating function $F(q,t)=\sum_{s\ge1}F_s(q)\,t^s$ satisfying the
\emph{validated} functional equation
\begin{equation}\label{eq:cat}
F(q,t)\;=\;E(t)\;+\;D(t)\,F(q,q^2t),
\qquad
D(t)=\frac{2q(q-1)\,t}{(1-qt)(1-q^2t)},
\quad
E(t)=\frac{2qt}{1-qt}\bigl(1+F(q,q)\bigr).
\end{equation}
Because the catalytic variable enters through the \emph{dilation} $t\mapsto q^2t$, equation
\eqref{eq:cat} is a linear $q$-difference (not polynomial-kernel / Bousquet-M\'elou--Jehanne) equation;
$F(q,q^2t)$ still depends on $t$. The series of interest is the diagonal specialization
\[
F(q):=F(q,1)\;=\;2q+2q^2+6q^3+2q^4+18q^5+6q^6+42q^7+18q^8+118q^9+\cdots\ \in\ \Z[[q]],
\]
the ``bulk-run'' block of the metric; the true geodesic series $U$ is a rational function of $F$ and its
companions. Telescoping \eqref{eq:cat} gives the closed form $F(q,t)=\sum_{k\ge0}\alpha_k(t)\,E(q^{2k}t)$
with $\alpha_k(t)=[2q(q-1)]^k\,t^k\,q^{k(k-1)}/(qt;q)_{2k}$ (verified symbolically), and at $t=1$,
\begin{equation}\label{eq:theta}
F(q)\;=\;2\bigl(1+F(q,q)\bigr)\,\Phi(q),
\qquad
\boxed{\;\Phi(q)\;=\;\sum_{k\ge0}\frac{2^k\,(q-1)^k\,q^{(k+1)^2}}{(q;q)_{2k+1}}\;}
\end{equation}
a partial-theta-type $q$-hypergeometric series. The exponent $q^{(k+1)^2}$ is a genuine theta/square
structure; the Euler-product denominator $(q;q)_{2k+1}=\prod_{i=1}^{2k+1}(1-q^i)$ spreads each theta
block densely across all degrees, so $F$ has \emph{no lacunary gaps}.

\section*{The question}

\begin{question}\label{q:main}
Is $F(q)=F(q,1)$ (equivalently $\Phi(q)$, equivalently $U(x)$) algebraic over $\Q(q)$, or transcendental?
\end{question}

\noindent\textbf{Characteristic-$p$ form (Christol).} For each prime $p$, $F\bmod p\in\F_p[[q]]$ is
algebraic over $\F_p(q)$ iff the sequence of its coefficients is $p$-automatic, iff its $p$-kernel is
finite. So Question~\ref{q:main} mod $p$ is the decidable-in-principle question: \emph{is the $p$-kernel
of $F\bmod p$ finite?}

\section*{What we know (evidence for transcendental / non-automatic)}

\begin{itemize}
\item \textbf{Maximal non-automaticity.} For $p=3$ the $p$-kernel sizes through level $k=3$ are
$1,4,13,39$, against the full ternary tree $1,4,13,40$ (deficit $1$); for $p=5$, $1,6,30$ against
$1,6,31$. Robust across overlap thresholds; the data ceiling is $k=3$ ($p=3$) / $k=2$ ($p=5$), since
$k=4$ needs $\sim 400$ terms. No eventual periodicity.
\item \textbf{Not rational.} Hankel determinants of $(u_n\bmod 3)$ are full-rank through size $60$;
Berlekamp--Massey linear complexity $\approx N/2$.
\item \textbf{No low-complexity char-$p$ Mahler relation.} Over $\F_3$, $F$ satisfies no relation
$\sum_{i=0}^n A_i(q)\,F(q^{3^i})=0$ of Frobenius-depth $\le4$ and $\deg A_i\le 32$ (nullspace search).
\item \textbf{The constituent blocks are unconditionally transcendental.} The four building-block
series ($\Sigma_1,\Sigma_0,S_1,S_0$) are integer series of radius $1$ with super-polynomial diagonal
growth, hence have $|q|=1$ as a natural boundary by P\'olya--Carlson, hence are transcendental over
$\Q(q)$. The difficulty is that $U$ is a \emph{ratio} of such blocks, in which the boundary
singularities can cancel.
\end{itemize}

\section*{What provably does \emph{not} work (and why)}

Every method below caps out against the same obstruction: the support of $F$ is \emph{dense}, and there
is an exact degree balance \#zeros $\sim$ \#poles $\sim 2N^2$ for the order-$N$ truncation remainder,
which is the algebraic signature of a natural boundary.
\begin{itemize}
\item \textbf{Degree / Pad\'e / Liouville / valuation methods} all reduce to comparing $2N^2$ against
$d\cdot P(N)$ with $P(N)=2N^2+2N+1$ \emph{exactly} (computed over $\Q$, $\F_3$, $\F_5$ identically):
the leading $2N^2$ terms match, ruling out degree $d=1$ (so $F$ is non-rational) but never reaching
$d\ge2$.
\item \textbf{Lacunarity / gap theorems} (Mahler-method transcendence, Adamczewski--Bell--Faverjon)
require \emph{sparse} support; the $(q;q)_{2k+1}$ denominator makes the support dense.
\item \textbf{Modular forms mod $p$:} $F$ is not a finite eta-quotient (log-derivative recovers a
growing set of $60+$ nonzero exponents).
\item \textbf{Parameterized (Frobenius) difference Galois theory:} in characteristic $p$,
$\varphi(f)=f^p$ \emph{exactly}, so the Frobenius orbit $\{f,f^p,\dots\}$ is purely algebraic over $f$
and $\varphi$-transcendence collapses to ordinary transcendence. The genuine $\sigma$-commuting
operator is the Euler derivation $\delta=t\partial_t$, but $\delta$ yields hyper-transcendence of
$F(q,t)$ in $t$, not algebraicity of the \emph{value} $F(q,1)$ over $\F_p(q)$ -- the wrong target.
\end{itemize}

\section*{The precise ask}

Is there a transcendence criterion that applies to the \emph{specialization at a regular point} ($t=1$,
where $D(1)$ is finite and nonzero) of a $\sigma$-transcendental solution of a linear $q$-difference
equation, when the $\sigma$-orbit of the specialization point accumulates at the singular fixed point
$t=0$? Equivalently: is the partial-theta series $\Phi(q)$ of \eqref{eq:theta} known, or amenable to a
known method (Rogers--Ramanujan / partial-theta transcendence, the work of Adamczewski--Dreyfus--Hardouin
and Dreyfus--Hardouin--Roques on functional relations for $q$-difference equations, or a $q\to1$
confluence argument that bridges the $q=0$-adic and $q=1$-adic behaviours)? Any pointer to the right
result, the right technique, or the right person would be very welcome.

\section*{A second, independent question (analytic route)}

The same transcendence conclusion follows by an entirely analytic route from one asymptotic
estimate for the \textbf{Hahn--Exton (third Jackson) $q$-Bessel function} $J^{(3)}_{3/2}(x;q)$ as
$q=e^{-\tau}\to1^-$. We have reduced that estimate, by exact identities, to a question about the
$q$-trigonometric pair of Koornwinder--Swarttouw. Sampled along $x_n=x_0q^{n/2}$, $J^{(3)}_\nu$
satisfies $u_{n+2}-B_nu_{n+1}+u_n=0$ with $B_n=(q^{\nu/2}+q^{-\nu/2})-q^{1-\nu/2}x_n^2$. Three exact,
verified facts:
\begin{itemize}
\item \emph{Half-integer ladder}: with $f_\nu={}_1\varphi_1(0;q^{\nu+1};q,qx^2)$,
$f_{3/2}(xq^{1/2})=\tfrac{(1-q^{1/2})(1-q^{3/2})}{q^{3/2}x^2}[f_{1/2}(x)-f_{-1/2}(x)]$ --- the exact
$q$-analogue of $j_1(X)=(\sin X/X-\cos X)/X$. So the question concerns only the $q$-sine/$q$-cosine
pair $\nu=\pm\tfrac12$.
\item \emph{Exact pair Casoratian}: on the $\nu=\tfrac12$ lattice, $u=J^{(3)}_{1/2}(x)$ and
$v=J^{(3)}_{-1/2}(xq^{-1/4})$ solve the same recurrence with constant Casoratian
$-\tfrac{(q^{1/2};q)_\infty(q^{3/2};q)_\infty}{(q;q)_\infty^2}q^{1/8}(q^{-1/4}-q^{1/4})$; the
amplitude \emph{product} of the pair is therefore pinned exactly, and only the phase offset
$\Delta$ (classically $\pi/2$) remains.
\item \emph{Soft target}: the application needs only $\Delta=\tfrac\pi2+O(\tau^{\delta})$
(equivalently, oscillatory-amplitude normalisation $1+O(\tau^\delta)$) for \emph{some} $\delta>0.$
\end{itemize}
In the variable $s=\log x$, the lattice step is $\tau/2$ while the classical Bessel turning-point
region has $s$-width $O(1)$: the recurrence is a \emph{fine} second-order discretisation of the
Bessel equation there, and the local model is classical Bessel (not $q$-Airy). The remaining step
looks like textbook asymptotic integration --- a discrete Levinson (Benzaid--Lutz) transport of the
recessive solution from $x=O(\tau)$ through the turning point, with per-step residuals $O(\tau^2)$.
\textbf{Question 2: has this uniform $q\to1^-$ (equivalently, coefficient-varying-with-stepsize)
Levinson analysis for the Hahn--Exton recurrence, or the amplitude normalisation of the
Koornwinder--Swarttouw $q$-trigonometric functions, been carried out anywhere in the literature?}
Zeros and orthogonality of these functions are well studied (Koelink--Swarttouw, Abreu, Bustoz,
Cardoso); we have not found the uniform confluent amplitude statement.

Either question closed makes the growth series $U$ unconditionally transcendental. Full code, data,
and the verified reductions are in the repository linked above.


\section{Sharpened form of the analytic question (July 2026)}

All series in the problem are now identified with Koornwinder--Swarttouw $q$-trigonometric
functions, $\cos(z;q^2)={}_1\phi_1(0;q;q^2,q^2z^2)$: with $z_0=\sqrt{2(1-q)}$, $Z=z_0/\sqrt q$,
$q=e^{-\tau}$,
\[
 1-\Sigma_1=\cos(Z;q^2),\qquad S_e=\cos(z_0;q^2),\qquad
 P_{12}=\tfrac{z_0}2\,\sin(z_0;q^2)-\cos(z_0;q^2),
\]
each an elementary identity (proved; verified to $60$ digits). The travel poles are the solutions of
$\cos(Z;q^2)=0$. The denominator bound $|S_e(q_m)|\ge0.35\sqrt{\tau_m}$ at every pole is proved by
elementary means (annulus representation, exact Gaussian moments, discrete Gronwall; no cited
asymptotic theorem). Exactly one analytic statement remains open, and it now has a one-sentence
form.

\medskip\noindent\textbf{Open question.} Let $A(\theta)=1/(re^{i\theta};q)_\infty$ with
$r=\sqrt{2(1-q)}\,e^{-\tau/2}$. Prove, with an explicit remainder uniform as $\tau\to0^+$,
\[
 \mathbb{E}\bigl[\operatorname{Re}A(\theta)\bigr]
 =\cos\bigl(z_0;q^2\bigr)\quad\text{for }\theta\sim N(\pi/2,\ \tau/2)
 \ \ \text{(this identity is proved)},
\]
evaluated to relative precision $O(\tau)$: the amplitude $A$ varies on the scale
$\sqrt{\tau/2}$, which is the Gaussian's own standard deviation, so no fixed-order Laplace
expansion converges relatively (measured: constant relative deficit $\approx0.39$ at every order
tested). Equivalently: prove the uniform confluent expansion of $\cos(z;q^2)$ at
$z=\sqrt{2(1-q)}$ to relative $O(\tau)$. This single matched-scale Gaussian convolution of an
explicit $q$-product is the entire remaining gap between the conditional and unconditional
transcendence of the series $U$; the corresponding relative correction is numerically
$\tfrac34\tau+O(\tau^2)$ at every computed pole, so the sharp constant appears to be
$\tfrac34$.

\end{document}
